\chapter*{ЗАКЛЮЧЕНИЕ}
\addcontentsline{toc}{chapter}{ЗАКЛЮЧЕНИЕ}

В рамках данной курсовой работы был разработан прототип системы бронирования автомобилей на выбранные даты. Система реализована в виде набора микросервисов, каждый из которых отвечает за собственную область: пользовательский интерфейс, авторизацию и учетные записи, автомобили, бронирования, платежи, статистику и агрегацию запросов.

Для backend-сервисов использован FastAPI, для пользовательского интерфейса --- React. В качестве СУБД используется PostgreSQL. Для асинхронной передачи событий и последующего построения статистики используется Apache Kafka. Авторизация реализована через собственный Identity Provider с поддержкой OpenID Connect, JWT и JWKs. Пользователи разделены на роли User и Admin, а доступ к административной статистике ограничен ролью Admin.

Система разворачивается в Kubernetes с использованием Helm charts. Для публикации наружу настроен Ingress Controller. Сборка Docker-образов, развертывание и запуск Newman-тестов выполняются через GitHub Actions.

Таким образом, были решены следующие задачи:

\begin{itemize}
  \item описана предметная область системы бронирования автомобилей;
  \item сформулированы функциональные требования к системе;
  \item спроектирована микросервисная архитектура;
  \item реализован пользовательский веб-интерфейс;
  \item реализован собственный Identity Provider с OpenID Connect;
  \item реализованы сервисы автомобилей, бронирования, платежей, статистики и gateway;
  \item настроена передача событий через Kafka;
  \item реализована token-based авторизация межсервисных запросов;
  \item подготовлено развертывание системы в Kubernetes через Helm;
  \item настроен CI/CD pipeline.
\end{itemize}

Цель курсовой работы достигнута.
